% unidades/16-experiencias.tex
\chapter{🌟 経験と回数 — Experiencias}
\unidadheader{16}{経験と回数}{Experiencias y Veces}{Contar experiencias y cuántas veces hiciste algo}

\begin{tcolorbox}[vocabbox, title={📌 Vocabulario Nuevo}]
\begin{tabularx}{\linewidth}{llX}
\toprule
\textbf{Japonés} & \textbf{Español} & \textbf{Tipo} \\
\midrule
\vocabitem{経験}{けいけん}{experiencia}{sust.}
\vocabitem{〜度}{ど}{vez / grado (contador)}{contador}
\vocabitem{〜回}{かい}{vez (contador de acciones)}{contador}
\vocabitem{初めて}{はじめて}{por primera vez}{adv.}
\vocabitem{一度も}{いちども}{ni una sola vez (+ neg.)}{adv.}
\vocabitem{まだ}{—}{todavía / aún (no)}{adv.}
\vocabitem{もう}{—}{ya (+ afirmativo)}{adv.}
\vocabitem{外国}{がいこく}{país extranjero}{sust.}
\vocabitem{海外}{かいがい}{el extranjero / exterior}{sust.}
\vocabitem{留学する}{りゅうがくする}{estudiar en el extranjero}{v. irr.}
\bottomrule
\end{tabularx}
\end{tcolorbox}

\subsection*{🗣️ Frases Clave}

\frase{\ruby{日本}{にほん}に\ruby{行}{い}ったことがありますか？}{¿Has ido a Japón alguna vez?}
\frase{はい、\ruby{二回}{にかい}\ruby{行}{い}ったことがあります。}
  {Sí, he ido dos veces.}
\frase{いいえ、まだ\ruby{行}{い}ったことがありません。}
  {No, todavía no he ido.}
\frase{\ruby{初}{はじ}めて\ruby{寿司}{すし}を\ruby{食}{た}べました。}
  {Comí sushi por primera vez.}

\subsection*{⚙️ Gramática}

\patron{Haber tenido experiencia}
  {V-た + ことがあります / ことがありません}
  {\ej{\ruby{富士山}{ふじさん}に\ruby{登}{のぼ}ったことがあります。}
    {He escalado el Monte Fuji.}
  \ej{スカイダイビングをしたことがありません。}
    {Nunca he hecho paracaidismo.}}

\patron{Cuántas veces: 〜回 / 〜度}
  {数字 + 回/度 + V}
  {\ej{\ruby{三回}{さんかい}\ruby{見}{み}たことがあります。}
    {Lo he visto tres veces.}
  \ej{一度も\ruby{遅刻}{ちこく}したことがありません。}
    {Nunca he llegado tarde.}}

\begin{tcolorbox}[ejerciciobox, title={📝 Ejercicios Rápidos}]
\begin{enumerate}
  \item Pregunta a alguien si alguna vez comió takoyaki (たこ焼き).
  \item Di algo que nunca has hecho usando ことがありません.
  \item Traduce: もう\ruby{宿題}{しゅくだい}をしましたか？
\end{enumerate}
\vspace{0.4em}
\textbf{Respuestas:}
1) たこ\ruby{焼}{や}きを\ruby{食}{た}べたことがありますか？\quad
2) (personal)\quad
3) ¿Ya hiciste la tarea?
\end{tcolorbox}

\begin{tcolorbox}[kanjibox, title={🀄 Kanjis de la Unidad}]
\begin{tabularx}{\linewidth}{cllXX}
\toprule
\textbf{Kanji} & \textbf{On} & \textbf{Kun} & \textbf{Significado} & \textbf{Vocab} \\
\midrule
\kanjirow{経}{けい}{へ}{pasar / experiencia}{\ruby{経験}{けいけん} / \ruby{経済}{けいざい}}
\kanjirow{験}{けん}{—}{experiencia / examen}{\ruby{経験}{けいけん} / \ruby{試験}{しけん}}
\kanjirow{度}{ど・たく}{たび}{vez / grado}{\ruby{一度}{いちど} / \ruby{温度}{おんど}}
\kanjirow{回}{かい}{まわ}{vuelta / vez}{\ruby{三回}{さんかい} / \ruby{回転}{かいてん}}
\kanjirow{初}{しょ}{はじ・うい}{primero / inicio}{\ruby{初}{はじ}めて / \ruby{初日}{しょにち}}
\bottomrule
\end{tabularx}
\end{tcolorbox}
